% Agitation des molécules d'un gaz à deux températures (L'essentiel, chapitre 18)
\begin{tikzpicture}[x=1cm,y=1cm]
\foreach \dx/\ech/\titre in {0/0.45/{gaz froid},4.6/0.9/{même gaz, plus chaud}}{
  \begin{scope}[shift={(\dx,0)}]
  \fill[pcEau] (0,0) rectangle (3.4,2.6);
  \draw[otInk!70,line width=0.8pt] (0,0) rectangle (3.4,2.6);
  \foreach \x/\y/\a in {0.5/0.5/30,1.3/1.9/200,2.2/0.7/120,2.9/1.8/250,0.8/1.4/-40,1.8/1.2/80,2.7/0.35/170,0.45/2.2/-10,1.55/0.4/150,2.45/2.25/300}{
    \fill[otNavy] (\x,\y) circle (0.07);
    \draw[otRed!80!black,line width=0.6pt,-{Stealth[length=1.3mm]}] (\x,\y) -- ++(\a:\ech);}
  \node[pcLegende,anchor=north] at (1.7,-0.05) {\titre};
  \end{scope}}
% choc sur la paroi
\draw[otMuted,line width=0.5pt,dash pattern=on 1.5pt off 1pt,-{Stealth[length=1.2mm]}] (7.35,1.15) -- (7.97,1.45) -- (7.35,1.75);
\draw[otGreenDark,line width=0.9pt,-{Stealth[length=1.6mm]}] (8.0,1.45) -- (8.55,1.45);
\node[pcLegende,anchor=west,text=otGreenDark,align=left] at (8.15,2.05) {force sur\\ la paroi};
\end{tikzpicture}
\par\smallskip\legende{Les molécules s'agitent en tous sens ; plus la température est élevée, plus elles sont rapides. Leurs chocs sur les parois exercent la force pressante.}
